\documentclass[11pt, a4paper]{article}

% ─── Packages ───────────────────────────────────────────────────────────────
\usepackage[T1]{fontenc}
\usepackage[utf8]{inputenc}

\usepackage[margin=1in]{geometry}
\usepackage{xcolor}
\usepackage{titlesec}
\usepackage{titling}
\usepackage{enumitem}
\usepackage{booktabs}
\usepackage{array}
\usepackage{tabularx}
\usepackage{mdframed}
\usepackage{hyperref}

\usepackage{parskip}
\usepackage{fancyhdr}
\usepackage{graphicx}
\usepackage{amsmath}
\usepackage{amssymb}
\usepackage{tcolorbox}
\tcbuselibrary{breakable, skins}

% ─── Colours ────────────────────────────────────────────────────────────────
\definecolor{primary}{HTML}{1A1A2E}      % deep navy
\definecolor{accent}{HTML}{E94560}       % vivid crimson-red
\definecolor{accentsoft}{HTML}{F5A623}   % amber highlight
\definecolor{subtle}{HTML}{4A4E69}       % muted slate
\definecolor{lightbg}{HTML}{F4F6FB}      % near-white blue tint
\definecolor{rulecol}{HTML}{E94560}

% ─── Hyperref ───────────────────────────────────────────────────────────────
\hypersetup{
  colorlinks=true,
  linkcolor=accent,
  urlcolor=accent,
  pdftitle={Adaptive Real-World Anomaly Intelligence System},
  pdfauthor={AI/ML Systems Research}
}

% ─── Header / Footer ────────────────────────────────────────────────────────
\pagestyle{fancy}
\fancyhf{}
\renewcommand{\headrulewidth}{0.4pt}
\renewcommand{\headrule}{\hbox to\headwidth{\color{accent}\leaders\hrule height \headrulewidth\hfill}}
\fancyhead[L]{\small\color{subtle}\textbf{ARAIS} \textbar\ Project Overview}
\fancyhead[R]{\small\color{subtle}AI/ML Systems Research}
\fancyfoot[C]{\small\color{subtle}\thepage}

% ─── Section Styling ────────────────────────────────────────────────────────
\titleformat{\section}
  {\large\bfseries\color{primary}}
  {\color{accent}\thesection.}{0.6em}{}[{\color{accent}\titlerule[0.6pt]}]

\titleformat{\subsection}
  {\normalsize\bfseries\color{subtle}}
  {\thesubsection.}{0.5em}{}

\titlespacing{\section}{0pt}{18pt}{8pt}
\titlespacing{\subsection}{0pt}{12pt}{4pt}

% ─── Custom Boxes ───────────────────────────────────────────────────────────
\newtcolorbox{gapbox}[2][]{
  breakable,
  colback=lightbg,
  colframe=accent,
  arc=3pt,
  boxrule=1pt,
  left=8pt, right=8pt, top=6pt, bottom=6pt,
  title={\small\bfseries\color{white}#2},
  fonttitle=\bfseries,
  coltitle=white,
  attach boxed title to top left={yshift=-2mm, xshift=4mm},
  boxed title style={colback=accent, arc=2pt, boxrule=0pt},
  #1
}

\newtcolorbox{visionbox}{
  breakable,
  colback=primary!4,
  colframe=primary,
  arc=3pt,
  boxrule=0.8pt,
  left=8pt, right=8pt, top=6pt, bottom=6pt,
}

\newtcolorbox{diffbox}{
  breakable,
  colback=accentsoft!10,
  colframe=accentsoft,
  arc=3pt,
  boxrule=0.8pt,
  left=8pt, right=8pt, top=6pt, bottom=6pt,
}

% ─── Utility commands ───────────────────────────────────────────────────────
\newcommand{\mytag}[1]{%
  \textcolor{white}{\colorbox{accent}{\footnotesize\bfseries\, #1\,}}%
}
\newcommand{\pillar}[1]{\textbf{\textcolor{primary}{#1}}}

% ════════════════════════════════════════════════════════════════════════════
\begin{document}
% ════════════════════════════════════════════════════════════════════════════

% ─── Title Block ────────────────────────────────────────────────────────────
\begin{center}
  {\fontsize{22}{26}\selectfont\bfseries\color{primary}
    Adaptive Real-World Anomaly\\[4pt]Intelligence System}\\[10pt]
  {\large\color{accent}\rule{6cm}{1.2pt}}\\[8pt]
  {\normalsize\color{subtle}
    \textit{Project Vision \textbar\ Research Overview \textbar\ System Architecture Philosophy}}\\[6pt]
  {\small\color{subtle} AI/ML Systems Research \textbf{·} April 2026}
\end{center}

\vspace{6pt}
\noindent\color{accent}\rule{\linewidth}{1.2pt}
\vspace{8pt}

% ─── Abstract strip ─────────────────────────────────────────────────────────
\begin{visionbox}
\noindent\textbf{\color{primary}Abstract.}\quad
Anomaly detection in production is a fundamentally different problem than anomaly detection in a research lab.
Real-world systems face noise, distributional shift, unseen failure modes, and hard latency budgets — 
conditions that standard benchmarks never stress-test.
\textbf{ARAIS} (\textit{Adaptive Real-World Anomaly Intelligence System}) addresses three under-solved 
structural gaps: unrealistic evaluation, rigid model trade-offs, and open-world detection failures.
It is designed not as a single model, but as an \textit{intelligent system} that evaluates, adapts, and detects
within a unified, operationally-aware framework.
\end{visionbox}

\vspace{4pt}

% ════════════════════════════════════════════════════════════════════════════
\section{Key Gaps in Current Systems}
% ════════════════════════════════════════════════════════════════════════════

Modern anomaly detection research has reached impressive benchmark numbers — 
yet production deployments routinely underperform, misfire, or silently degrade.
The disconnect is not a model quality problem; it is a \textit{systems thinking} problem.
Three structural gaps account for the majority of this production gap.

\vspace{8pt}

% ── Gap 1 ───────────────────────────────────────────────────────────────────
\begin{gapbox}{\,Gap 1\; —\; Unrealistic Evaluation}
\textbf{Why it exists.}\;
Academic benchmarks optimise for reproducibility: static, pre-labelled, class-balanced datasets where 
the test distribution matches the training distribution. This design choice is reasonable for controlled 
science but catastrophically wrong as a proxy for production.

\smallskip
\textbf{Why it is critical.}\;
Deployed systems face sensor drift, mislabelled telemetry, temporally-lagged anomaly annotations, 
and rolling concept drift. A model that scores 0.97 AUC on a frozen dataset may produce 40\,\% 
false-positive rates under mild real-world noise — a failure mode that is completely invisible 
on the benchmark.

\smallskip
\textbf{How it limits production ML.}\;
Teams ship models with inflated confidence, then burn engineering cycles on post-deployment 
fire-fighting. Evaluation becomes a compliance ritual rather than a predictive signal of production 
health. Without temporal and noise-aware evaluation harnesses, model selection and iteration 
are fundamentally misguided.
\end{gapbox}

\vspace{8pt}

% ── Gap 2 ───────────────────────────────────────────────────────────────────
\begin{gapbox}{\,Gap 2\; —\; Model Trade-off Limitations}
\textbf{Why it exists.}\;
The dominant paradigm in anomaly detection is to train a single model — or an ensemble — 
and deploy it as a static artefact. Model selection happens offline, optimised for a single 
objective (usually recall or AUC), and then frozen for the lifecycle of the deployment.

\smallskip
\textbf{Why it is critical.}\;
Production environments are not stationary. A microservice anomaly detector may face 10$\times$ 
traffic spikes, a GPU-constrained edge device, a high-risk regulatory window requiring higher 
precision, or a degraded upstream pipeline that changes input statistics.
No single model is simultaneously optimal for high-load, low-latency, high-accuracy, 
and low-compute scenarios.

\smallskip
\textbf{How it limits production ML.}\;
Fixed models create brittle SLOs. Operators either over-provision compute 
(to cover worst-case latency) or accept accuracy degradation under load. 
The system has no mechanism to reason about its own performance-cost trade-off 
or escalate decisions based on contextual risk — it simply fires the same detector 
regardless of runtime state.
\end{gapbox}

\vspace{8pt}

% ── Gap 6 ───────────────────────────────────────────────────────────────────
\begin{gapbox}{\,Gap 6\; —\; Open-World Detection Failure}
\textbf{Why it exists.}\;
Supervised and semi-supervised anomaly detectors are trained on known anomaly classes or on 
a closed normal distribution. The implicit assumption — that the test distribution is a subset 
of the training distribution — is violated the moment a genuinely novel fault mode appears.

\smallskip
\textbf{Why it is critical.}\;
Novel anomalies are precisely the events that matter most: zero-day system intrusions, 
previously-unseen equipment failure signatures, emergent data pipeline corruptions.
Systems that cannot handle out-of-distribution (OOD) inputs either silently misclassify them 
as normal or force them into the nearest known class — both outcomes are operationally dangerous.

\smallskip
\textbf{How it limits production ML.}\;
Overconfident closed-world detectors produce falsely high softmax probabilities for OOD inputs,
masking the very events they should be flagging. The absence of uncertainty quantification means 
there is no signal to trigger human review, escalation, or model update cycles. 
The system does not know what it does not know.
\end{gapbox}

\vspace{2pt}

% ════════════════════════════════════════════════════════════════════════════
\section{Vision of the Solution}
% ════════════════════════════════════════════════════════════════════════════

ARAIS is designed around a single unifying principle:
\textit{a production anomaly detection system must be as intelligent about its own operating conditions 
as it is about the data it monitors.}

The system vision spans four interconnected capabilities:

\vspace{8pt}
\noindent\pillar{2.1\quad Realistic, Dynamic Evaluation Infrastructure}
\vspace{4pt}

\begin{visionbox}
Standard evaluation is replaced with a \textit{dynamic evaluation harness} that stress-tests models 
under conditions that mirror real deployments: injected sensor noise, synthetic concept drift 
at configurable rates, temporally delayed anomaly labels, class imbalance schedules, 
and operational constraints such as throughput budgets and memory ceilings.

The harness produces not a single performance number but a \textit{performance profile} — 
a multi-dimensional characterisation of model behaviour across noise levels, drift intensities, 
and latency regimes. This profile drives model selection and informs system-level decisions 
throughout the deployment lifecycle.

The goal is to make evaluation a \textit{continuous, live process} rather than a 
pre-deployment gate: models are tracked, re-evaluated against incoming distribution statistics, 
and flagged when their profile diverges from deployment conditions.
\end{visionbox}

\vspace{10pt}
\noindent\pillar{2.2\quad Runtime-Adaptive Model Selection and Orchestration}
\vspace{4pt}

\begin{visionbox}
A static model cannot reason about its own operational context.
ARAIS introduces a \textit{model orchestration layer} that continuously monitors runtime signals —
inference latency, system load, input distribution statistics, and risk tier of the 
current decision window — and dynamically routes inference to the most appropriate detector.

This is not simple model ensembling. It is contextual orchestration: a lightweight, fast detector 
handles high-volume, low-risk windows; a higher-capacity model is invoked for anomalies flagged 
as high-risk or during low-load periods where latency budgets are relaxed.
The orchestrator continuously re-learns the cost-accuracy frontier of the available model pool, 
adapting as new models are introduced or as system conditions change.

The outcome is a system where accuracy, latency, and compute are treated as 
\textit{jointly-optimised runtime objectives} rather than fixed offline trade-offs.
\end{visionbox}

\vspace{10pt}
\noindent\pillar{2.3\quad Open-World Detection with Calibrated Uncertainty}
\vspace{4pt}

\begin{visionbox}
ARAIS treats anomaly detection as an open-world inference problem from the ground up.
Every inference path produces not just a detection decision but a \textit{structured uncertainty signal}:
epistemic uncertainty over model parameters, aleatoric uncertainty over input ambiguity, 
and a dedicated OOD proximity score that quantifies how far the current input lies from 
the training distribution.

Inputs that exceed OOD thresholds are not forced into known categories.
They are routed to a \textit{novel-anomaly triage pipeline}: flagged for human review, 
logged to an active-learning queue for future retraining, 
and optionally handled by a conservative fallback that errs toward escalation over silence.

This architecture means the system has a principled answer to the question:
\textit{"Is this anomaly something I recognise, or something genuinely new?"} —
a question that current closed-world systems cannot even pose.
\end{visionbox}

\vspace{10pt}
\noindent\pillar{2.4\quad End-to-End Intelligent Integration}
\vspace{4pt}

\begin{visionbox}
The three capabilities above are not independent modules — they are feedback-coupled subsystems 
within a single operational loop. Evaluation outputs feed model selection;
model selection decisions inform which evaluation dimensions to prioritise;
OOD signals trigger re-evaluation and orchestration updates.

This closed-loop architecture means the system actively maintains its own 
operational health rather than passively serving requests.
It is anomaly detection designed to function as a \textit{first-class production service}, 
not a research model wrapped in an API endpoint.
\end{visionbox}

% ════════════════════════════════════════════════════════════════════════════
\section{What Makes It Different}
% ════════════════════════════════════════════════════════════════════════════

The research and industry landscape contains many capable anomaly detectors.
ARAIS is differentiated not by the sophistication of any individual model 
but by a set of \textit{system-level design commitments} that are largely absent from existing work.

\vspace{8pt}

\begin{diffbox}
\noindent\mytag{SHIFT 1}\quad\textbf{From Model-Centric to System-Level Intelligence}

\smallskip
Most anomaly detection research optimises a single model on a fixed task.
ARAIS optimises an \textit{operational system} across a dynamic task space.
The unit of intelligence is not the model; it is the system's ability to select, 
compose, evaluate, and replace models in response to changing conditions.
This mirrors how mature software systems handle reliability: through architecture and 
redundancy, not by betting on a single perfect component.
\end{diffbox}

\vspace{8pt}

\begin{diffbox}
\noindent\mytag{SHIFT 2}\quad\textbf{Deployment Realism as a First-Class Design Constraint}

\smallskip
Every design decision in ARAIS is stress-tested against the question:
\textit{"Does this work in a real deployment, not just in a notebook?"}
Latency budgets, memory ceilings, data quality degradation, 
and operational cost are not afterthoughts — they are embedded into the evaluation harness, 
the orchestration logic, and the detection pipeline from day one.
This is a deliberate departure from accuracy-first research culture 
toward a \textit{deployment-first engineering culture}.
\end{diffbox}

\vspace{8pt}

\begin{diffbox}
\noindent\mytag{SHIFT 3}\quad\textbf{Self-Aware Architecture with Uncertainty as a Core Signal}

\smallskip
ARAIS is designed to know the boundaries of its own competence.
Uncertainty quantification is not a metric reported at evaluation time; 
it is a live signal that shapes inference routing, escalation decisions, 
and human-in-the-loop triggers.
A system that can communicate \textit{"I am not confident about this input"} 
is qualitatively safer than one that silently outputs a high-confidence wrong answer.
This self-awareness is the foundational property that makes ARAIS trustworthy in high-stakes domains —
not just accurate under benign conditions.
\end{diffbox}

\vspace{8pt}

\begin{diffbox}
\noindent\mytag{SHIFT 4}\quad\textbf{Unified Integration of Evaluation, Adaptation, and Detection}

\smallskip
Existing systems treat evaluation, model selection, and detection as three separate concerns, 
often owned by three separate teams. ARAIS unifies them into a single operational loop 
where each component continuously informs the others.

\smallskip
\noindent The resulting architecture is summarised below:

\vspace{6pt}
\begin{center}
\renewcommand{\arraystretch}{1.35}
\begin{tabularx}{0.94\linewidth}{>{\bfseries\color{primary}}l X}
\toprule
\color{primary}Component & \color{primary}Function \\
\midrule
Evaluation Harness & Continuous, noise-aware, drift-injecting model profiling \\
Orchestration Layer & Runtime-adaptive routing across cost-accuracy frontier \\
OOD \& Uncertainty Engine & Principled handling of novel and ambiguous inputs \\
Feedback Loop & Closed-loop coupling across all three components \\
\bottomrule
\end{tabularx}
\end{center}

\vspace{4pt}
No existing open-source or commercial system integrates all four components 
into a coherent operational architecture.
ARAIS fills that gap.
\end{diffbox}

% ════════════════════════════════════════════════════════════════════════════
\section*{Closing Statement}
% ════════════════════════════════════════════════════════════════════════════
\addcontentsline{toc}{section}{Closing Statement}
\vspace{-4pt}
\noindent\color{accent}\rule{\linewidth}{0.6pt}
\vspace{6pt}

\begin{visionbox}
The next frontier in anomaly detection is not a better model architecture.
It is \textit{better systems thinking} — the ability to evaluate honestly, 
adapt intelligently, and detect openly in the full complexity of real-world deployments.

ARAIS represents a research commitment to building anomaly detection 
that is not merely performant under idealised conditions, but \textbf{trustworthy under operational ones}.

It is anomaly intelligence designed for the world as it is, not as benchmarks imagine it to be.
\end{visionbox}

\vspace{10pt}
\noindent\color{accent}\rule{\linewidth}{1.2pt}
\begin{center}
  \small\color{subtle}
  \textit{Adaptive Real-World Anomaly Intelligence System — Research Vision Document}\\
  \textit{For distribution to senior engineers, ML practitioners, and technical leadership.}
\end{center}

\end{document}